% Manual de usuario de DevStudio
% Compilar: /usr/bin/latexbuild manual-usuario.tex
% Preámbulo compartido por el reporte, los manuales y el contrato.
% Se compila con XeLaTeX mediante /usr/bin/latexbuild.
\documentclass[11pt,a4paper]{article}

\usepackage{fontspec}
\setmonofont{JetBrainsMono Nerd Font Mono}[Scale=0.82]
\usepackage[spanish,es-nodecimaldot,es-noshorthands]{babel}

\usepackage[top=2.4cm,bottom=2.4cm,left=2.4cm,right=2.4cm]{geometry}
\usepackage{fancyhdr}
\usepackage{titlesec}
\usepackage{xcolor}
\usepackage{booktabs}
\usepackage{tabularx}
\usepackage{array}
\usepackage{enumitem}
\usepackage{float}
\usepackage{graphicx}
\usepackage{tikz}
\usetikzlibrary{positioning,arrows.meta}
\usepackage{listings}
\usepackage{xurl}

\makeatletter
\renewcommand*\l@section[2]{%
  \ifnum \c@tocdepth >\z@
    \addpenalty\@secpenalty
    \addvspace{1.0em \@plus\p@}%
    \setlength\@tempdima{2.6em}%
    \begingroup
      \parindent \z@ \rightskip \@pnumwidth
      \parfillskip -\@pnumwidth
      \leavevmode \bfseries
      \advance\leftskip\@tempdima
      \hskip -\leftskip
      #1\nobreak\hfil \nobreak\hb@xt@\@pnumwidth{\hss #2}\par
    \endgroup
  \fi}
\makeatother
\usepackage[hidelinks]{hyperref}

\definecolor{tnBlue}{HTML}{0EA5E9}
\definecolor{tnDark}{HTML}{0F172A}
\definecolor{codeGray}{HTML}{F4F6F8}
\definecolor{codeGreen}{HTML}{2E7D32}
\definecolor{codeRed}{HTML}{C62828}

\newcolumntype{Y}{>{\raggedright\arraybackslash}X}

\lstset{
  basicstyle=\ttfamily\footnotesize,
  breaklines=true,
  breakatwhitespace=false,
  frame=tb,
  framerule=0.4pt,
  rulecolor=\color{tnDark!40},
  framesep=5pt,
  keywordstyle=\color{tnBlue}\bfseries,
  commentstyle=\color{codeGreen}\itshape,
  stringstyle=\color{codeRed},
  showstringspaces=false,
  tabsize=2,
  columns=fullflexible,
  keepspaces=true,
  extendedchars=true,
  captionpos=t,
  xleftmargin=0pt,
  xrightmargin=0pt,
  aboveskip=0.9em,
  belowskip=0.9em,
  literate=
    {á}{{\'a}}1 {é}{{\'e}}1 {í}{{\'i}}1 {ó}{{\'o}}1 {ú}{{\'u}}1
    {Á}{{\'A}}1 {É}{{\'E}}1 {Í}{{\'I}}1 {Ó}{{\'O}}1 {Ú}{{\'U}}1
    {ñ}{{\~n}}1 {Ñ}{{\~N}}1 {ü}{{\"u}}1 {Ü}{{\"U}}1
    {¿}{{?`}}1 {¡}{{!`}}1
    {·}{{\textperiodcentered}}1
}

\pagestyle{fancy}
\fancyhf{}
\fancyhead[L]{\small Alfredo Ramírez Candia}
\fancyhead[R]{\small Diseño de Aplicaciones Web}
\fancyfoot[C]{\thepage}
\renewcommand{\headrulewidth}{0.4pt}

\titleformat{\section}
  {\Large\bfseries\color{tnDark}}{\thesection.}{0.5em}{}
\titleformat{name=\section,numberless}
  {\Large\bfseries\color{tnDark}}{}{0em}{}
\titleformat{\subsection}
  {\large\bfseries\color{tnBlue}}{\thesubsection}{0.5em}{}
\titlespacing*{\section}{0pt}{1.4em}{0.6em}
\titlespacing*{\subsection}{0pt}{1.05em}{0.35em}

\addto\captionsspanish{%
  \renewcommand{\tablename}{Tabla}%
  \renewcommand{\figurename}{Figura}%
  \renewcommand{\lstlistingname}{Listado}%
}

\newcommand{\datosPdf}[1]{%
  \hypersetup{%
    pdftitle={#1},%
    pdfauthor={Alfredo Ramírez Candia},%
    pdflang={es}%
  }%
}

\newcommand{\portada}[3]{%
  \begin{titlepage}
    \centering
    \vspace*{0.8cm}
    {\large\bfseries BENEMÉRITA UNIVERSIDAD AUTÓNOMA DE PUEBLA\par}
    \vspace{0.28cm}
    {\large FACULTAD DE CIENCIAS DE LA COMPUTACIÓN\par}
    \vspace{0.28cm}
    {\large LICENCIATURA/INGENIERÍA EN CIENCIAS DE LA COMPUTACIÓN\par}
    \vspace{1.35cm}
    \rule{\linewidth}{0.5pt}\par
    \vspace{0.45cm}
    {\Huge\bfseries #1\par}
    \vspace{0.38cm}
    {\Large #2\par}
    \vspace{0.32cm}
    {\large Asignatura: Diseño de Aplicaciones Web\par}
    \vspace{0.45cm}
    \rule{\linewidth}{0.5pt}\par
    \vspace{1.5cm}
    \begin{tabular}{rl}
      \textbf{Docente:}   & Miguel Ángel Vargas Lomelí \\[0.55cm]
      \textbf{Alumno:}    & Alfredo Ramírez Candia \\[0.55cm]
      \textbf{Matrícula:} & 202375947 \\[0.55cm]
      \textbf{Fecha:}     & 22 de noviembre de 2026 \\
    \end{tabular}
    \vfill
    {\small #3\par}
  \end{titlepage}
  \setcounter{page}{2}
}

\setcounter{tocdepth}{2}
\setcounter{secnumdepth}{2}
\setlist[enumerate]{leftmargin=1.5em,itemsep=0.45em,topsep=0.4em}
\setlist[itemize]{leftmargin=1.4em,itemsep=0.28em,topsep=0.35em}
\emergencystretch=2em
\urlstyle{tt}
\renewcommand{\arraystretch}{1.18}

\datosPdf{Manual de usuario. DevStudio}

\begin{document}
\portada{Manual de usuario}{Sitio, portal y panel de DevStudio}{Manual de usuario}

\tableofcontents
\clearpage

\section{Para quién es este manual}

Este manual explica cómo usar DevStudio sin entrar al código. El sitio público
sirve para conocer el estudio. El portal es de quien ya es cliente. El panel
es de quien administra el contenido y la bandeja.

En la demostración local el sitio está en \url{http://localhost:8080}. Hace
falta un navegador actual, con JavaScript, y conexión si el servidor no está
en la misma máquina.

\section{Entrar y salir}

\subsection{Crear cuenta}

\begin{enumerate}
  \item Abrir \path{/pages/registro.html}.
  \item Escribir nombre, correo y empresa. El teléfono es opcional.
  \item Escribir la contraseña dos veces.
  \item Pulsar \textbf{Crear cuenta}.
\end{enumerate}

La contraseña tiene entre 8 y 200 caracteres, con al menos una mayúscula, una
minúscula y un número. Si el correo ya existe, hay que iniciar sesión o usar
otro. Si las dos contraseñas no coinciden, el formulario lo indica y conserva
el resto.

\subsection{Iniciar sesión}

En \path{/pages/login.html} se capturan correo y contraseña y se pulsa
\textbf{Entrar}. Un cliente llega al resumen del portal. Un administrador
llega al tablero. Tras cinco intentos fallidos con el mismo correo o desde la
misma dirección, el acceso espera 15 minutos.

\subsection{Cerrar sesión}

En el portal, el botón \textbf{Cerrar sesión} está en Resumen, debajo de la
actividad reciente. También está en el menú del monograma, en cualquier
página. Al salir se vuelve al inicio y el encabezado muestra de nuevo
\textbf{Iniciar sesión}.

\subsection{Cuentas de la demostración}

\begin{table}[H]
  \centering
  \caption{Cuentas para recorrer la demostración local.}
  \begin{tabular}{lll}
    \toprule
    Uso & Correo & Contraseña \\
    \midrule
    Administración & \texttt{admin@devstudio.local} & \texttt{Admin\#2026} \\
    Cliente, Laura Méndez & \texttt{cliente@devstudio.local} & \texttt{Cliente\#2026} \\
    \bottomrule
  \end{tabular}
\end{table}

\section{Sitio público}

El menú fijo tiene Soluciones, Casos, Cómo trabajamos y Contacto. A la derecha
están el cambio de día y noche, la cuenta y el botón \textbf{Hablar de mi
negocio}.

\begin{itemize}
  \item \textbf{Inicio.} Presenta el estudio y lleva al diagnóstico.
  \item \textbf{Soluciones.} Lista los servicios activos. Cada uno abre su detalle.
  \item \textbf{Casos.} Muestra proyectos publicados. Se puede filtrar por tecnología.
  \item \textbf{Cómo trabajamos.} Explica el modo de trabajo del estudio.
  \item \textbf{Contacto.} Pide nombre, correo, asunto y mensaje.
  \item \textbf{Diagnóstico.} En \path{/pages/cotizacion.html} se describe el negocio sin crear cuenta.
\end{itemize}

Si faltan campos, el formulario señala el error y no borra lo ya escrito. Si
se envía demasiadas veces en poco tiempo, hay que esperar.

El tema elegido se recuerda en ese navegador. Crear cuenta no está en el menú
principal: el enlace está en la página de inicio de sesión.

\section{Portal del cliente}

Solo entra quien tiene rol de cliente. Otro rol que abra estas páginas vuelve
al sitio. Las secciones son Resumen, Proyectos, Solicitudes, Calendario y
Perfil.

\subsection{Resumen}

Muestra el nombre, la empresa y cuatro cifras: empresa, proyectos activos,
solicitudes abiertas y proyectos en total. Debajo está la actividad reciente.
La foto abre Perfil. El botón de cerrar sesión está al final de esta página.

\subsection{Proyectos y avance}

Mis proyectos lista los trabajos de la empresa de la cuenta. Cada fila tiene
una barra de 0 a 100. El detalle repite la barra y el historial: fecha y nota
de cada actualización. El cliente solo consulta. No hay un campo para cambiar
el porcentaje.

\subsection{Solicitudes}

La lista muestra las solicitudes de esa cuenta y su estado. El detalle incluye
la respuesta cuando el estudio ya la escribió. Desde el portal también se puede
crear una solicitud nueva. No se ven las de otras empresas.

\subsection{Calendario}

La página muestra el mes, con botones para el mes anterior y el siguiente. Un
día con reuniones indica cuántas hay. Al elegirlo se ven el título, el horario,
el proyecto, el estado y, si el estudio ya los escribió, la respuesta y los
puntos.

Para agendar se eligen proyecto, título, inicio y fin. La reunión aparece al
momento con estado \textbf{Solicitada}. La duración válida va de 15 minutos a
4 horas y no puede cruzarse con otra reunión activa de la misma cuenta. El
horario se interpreta en la hora de la Ciudad de México.

\textbf{Cancelar reunión} solo aparece si la cita sigue solicitada o
confirmada y todavía no empieza. Una cita cancelada permanece en el día, con
la etiqueta correspondiente, y ya no se puede cancelar otra vez.

\subsection{Perfil y foto}

Perfil muestra cómo ve el equipo a la persona: nombre, cargo, empresa, frase
y ciudad. El correo no se edita ahí. \textbf{Guardar perfil} actualiza el
resto. Más abajo se cambia la contraseña; pide la actual y la nueva dos veces.

La foto es el círculo grande. Al pasar el cursor aparece un lápiz. Al pulsarlo
se elige una imagen JPG, PNG o WebP de hasta 2 MB. La misma foto se ve en el
menú de la cuenta. Si todavía no hay foto, se muestran las iniciales.

\section{Panel de administración}

Entra con la cuenta de administración. El tablero resume conteos y actividad.
Desde el menú interno se abren:

\begin{itemize}
  \item usuarios, su rol y si la cuenta está activa;
  \item clientes y sus datos de empresa;
  \item servicios y tecnologías del catálogo público;
  \item proyectos, incluida la publicación en casos;
  \item solicitudes, su estado y la conversión a proyecto;
  \item mensajes de contacto.
\end{itemize}

Confirmar una reunión o cambiar el porcentaje de avance no tiene pantalla en
este panel. Esas dos acciones se hacen con la API, como explica el manual
técnico. Un cliente no puede usarlas.

\section{Preguntas frecuentes}

\begin{itemize}
  \item \textbf{No puedo entrar.} Revisar mayúsculas y, si hubo varios fallos, esperar 15 minutos.
  \item \textbf{No veo los proyectos de otra empresa.} Es el comportamiento esperado.
  \item \textbf{La reunión no se agenda.} El fin debe ir después del inicio, durar entre 15 minutos y 4 horas, y no empalmarse con otra cita activa.
  \item \textbf{La foto no cambia.} El archivo tiene que ser JPG, PNG o WebP y pesar como máximo 2 MB.
  \item \textbf{El día y la noche vuelven al otro modo.} El navegador recuerda la última elección en ese equipo.
\end{itemize}

\section{Soporte}

En la demostración de la materia el contacto visible del sitio es
\texttt{contacto@devstudio.local}. Para una instalación propia, quien opera el
servidor es quien restablece una contraseña olvidada: el sitio no envía correo
de recuperación.

\end{document}
